\section{Modello di Navigazione}

Il modello di navigazione e' stato ricavato dalla struttura reale di \texttt{src/app}. In Next.js App Router ogni cartella che contiene \texttt{page.tsx} diventa una pagina visitabile; per questo motivo il diagramma seguente non descrive funzionalita' ipotetiche, ma solo route presenti nel codice.

\subsection{Legenda}

\begin{table}[H]
\centering
\renewcommand{\arraystretch}{1.25}
\begin{tabular}{|c|p{10cm}|}
\hline
\textbf{Sigla} & \textbf{Significato} \\
\hline
VP & Vista pubblica, accessibile senza autenticazione. \\
\hline
VF & Vista di form, usata per login o registrazione. \\
\hline
VD & Vista dashboard o vista di riepilogo. \\
\hline
VG & Vista gestionale, usata per modificare dati o configurazioni. \\
\hline
VA & Vista assistita per il paziente, con percorso semplice e pochi pulsanti. \\
\hline
API & Route non navigabile come pagina, usata dal codice applicativo. \\
\hline
\end{tabular}
\caption{Legenda del modello di navigazione}
\label{tab:legenda-navigazione}
\end{table}

\subsection{Albero di navigazione reale}

\begin{figure}[H]
\centering
\resizebox{0.98\textwidth}{!}{
\begin{tikzpicture}[
    route/.style={rectangle, draw=aliceblue, fill=alicelight, thick, rounded corners=3pt, text width=3.7cm, minimum height=0.85cm, align=center, font=\scriptsize},
    role/.style={rectangle, draw=alicegreen, fill=alicegreen!10, thick, rounded corners=3pt, text width=3.4cm, minimum height=0.85cm, align=center, font=\scriptsize\bfseries},
    api/.style={rectangle, draw=aliceorange, fill=aliceorange!10, thick, rounded corners=3pt, text width=4cm, minimum height=0.85cm, align=center, font=\scriptsize},
    line/.style={thick, color=aliceblue, -{Stealth[length=2.3mm]}},
    apiLine/.style={thick, dashed, color=aliceorange, -{Stealth[length=2.3mm]}},
]

\node[route] (home) at (0, 0) {\texttt{/}\\Landing\\VP};

\node[role] (opArea) at (-5.2, -1.55) {Area Operatore};
\node[role] (pzArea) at (5.2, -1.55) {Area Paziente};
\draw[line] (home) -- (opArea);
\draw[line] (home) -- (pzArea);

\node[route] (opLogin) at (-8.3, -3.2) {\texttt{/operatore/login}\\Login\\VF};
\node[route] (opReg) at (-5.2, -3.2) {\texttt{/operatore/registrati}\\Registrazione\\VF};
\node[route] (opDash) at (-2.1, -3.2) {\texttt{/operatore/dashboard}\\Dashboard\\VD};
\draw[line] (opArea) -- (opLogin);
\draw[line] (opArea) -- (opReg);
\draw[line] (opArea) -- (opDash);

\node[route] (opDetail) at (-2.1, -4.9) {\texttt{/operatore/dashboard/pazienti/[id]}\\Dettaglio paziente\\VD};
\node[route] (opManage) at (-2.1, -6.6) {\texttt{?tab=gestisci}\\Gestione dati, farmaci, esercizi\\VG};
\node[api] (aiApi) at (-6.7, -5.75) {\texttt{/api/ai/suggerimento-operatore}\\Suggerimento AI\\API};
\draw[line] (opDash) -- (opDetail);
\draw[line] (opDetail) -- (opManage);
\draw[apiLine] (opDetail) -- (aiApi);

\node[route] (pzLogin) at (2.1, -3.2) {\texttt{/paziente/login}\\Login\\VF};
\node[route] (pzReg) at (5.2, -3.2) {\texttt{/paziente/registrati}\\Registrazione\\VF};
\node[route] (pzHome) at (8.3, -3.2) {\texttt{/paziente/home}\\Home\\VA};
\draw[line] (pzArea) -- (pzLogin);
\draw[line] (pzArea) -- (pzReg);
\draw[line] (pzArea) -- (pzHome);

\node[route] (pzMood) at (4.3, -4.9) {\texttt{/paziente/umore}\\Umore\\VA};
\node[route] (pzEx) at (7.4, -4.9) {\texttt{/paziente/esercizi}\\Esercizi\\VA};
\node[route] (pzMed) at (10.5, -4.9) {\texttt{/paziente/farmaci}\\Farmaci\\VA};
\draw[line] (pzHome) -- (pzMood);
\draw[line] (pzHome) -- (pzEx);
\draw[line] (pzHome) -- (pzMed);

\end{tikzpicture}
}
\caption{Albero di navigazione aggiornato alla struttura reale delle pagine}
\label{fig:navigazione-reale}
\end{figure}

\subsection{Descrizione delle viste}

\subsubsection{Area comune}
\begin{itemize}
    \item \textbf{\texttt{/}:} pagina iniziale. Presenta l'accesso alle due aree del sistema. Se l'utente e' gia' autenticato, il file \texttt{proxy.ts} lo reindirizza alla home corretta in base al ruolo.
\end{itemize}

\subsubsection{Area operatore}
\begin{itemize}
    \item \textbf{\texttt{/operatore/login}:} form di accesso per l'operatore sanitario.
    \item \textbf{\texttt{/operatore/registrati}:} form di registrazione dell'operatore.
    \item \textbf{\texttt{/operatore/dashboard}:} vista principale dell'operatore con elenco pazienti, ricerca, filtri e metriche sintetiche.
    \item \textbf{\texttt{/operatore/dashboard/pazienti/[id]}:} dettaglio del paziente con dati anagrafici, farmaci, esercizi, umore, storico settimanale e suggerimento AI.
    \item \textbf{\texttt{?tab=gestisci}:} stato della pagina di dettaglio, non nuova route. Serve a modificare dati, farmaci ed esercizi.
\end{itemize}

\subsubsection{Area paziente}
\begin{itemize}
    \item \textbf{\texttt{/paziente/login}:} form di accesso del paziente.
    \item \textbf{\texttt{/paziente/registrati}:} form di registrazione del paziente.
    \item \textbf{\texttt{/paziente/home}:} home semplificata con tre scelte principali.
    \item \textbf{\texttt{/paziente/umore}:} registrazione o aggiornamento dell'umore giornaliero.
    \item \textbf{\texttt{/paziente/esercizi}:} elenco degli esercizi assegnati e percorso guidato step-by-step.
    \item \textbf{\texttt{/paziente/farmaci}:} elenco dei farmaci del giorno, divisi per fascia oraria, con conferma di assunzione.
\end{itemize}

\subsection{Regole di protezione e reindirizzamento}

La navigazione privata viene controllata dal file \texttt{src/proxy.ts}. Il proxy legge la sessione Supabase e applica poche regole:

\begin{itemize}
    \item utenti non autenticati verso \texttt{/operatore/dashboard} vengono mandati a \texttt{/operatore/login};
    \item utenti non autenticati verso pagine private \texttt{/paziente/} vengono mandati a \texttt{/paziente/login};
    \item utenti autenticati sulla landing vengono mandati alla home del proprio ruolo;
    \item utenti autenticati sulle pagine pubbliche di login o registrazione vengono reindirizzati alla propria area.
\end{itemize}

\subsection{Principi di navigazione}

\begin{enumerate}
    \item \textbf{Separazione per ruolo:} operatore e paziente hanno percorsi distinti.
    \item \textbf{Percorso breve per il paziente:} dalla home paziente si arriva a umore, esercizi e farmaci con un solo click.
    \item \textbf{Dashboard unica per l'operatore:} l'operatore parte dalla dashboard e poi entra nel dettaglio del paziente.
    \item \textbf{Nessuna pagina non implementata:} il modello contiene solo route realmente presenti nel codice.
\end{enumerate}

\newpage
